\documentclass[11pt]{article}

\usepackage[margin=1in]{geometry}
\usepackage{amsmath,amssymb,amsthm,mathtools}
\usepackage{enumitem}
\usepackage[T1]{fontenc}
\usepackage{lmodern}
\usepackage{microtype}

\newtheorem{theorem}{Theorem}
\newtheorem{lemma}[theorem]{Lemma}
\newtheorem{proposition}[theorem]{Proposition}
\newtheorem{corollary}[theorem]{Corollary}
\theoremstyle{definition}
\newtheorem{definition}[theorem]{Definition}
\theoremstyle{remark}
\newtheorem{remark}[theorem]{Remark}

\newcommand{\area}{\operatorname{area}}
\newcommand{\dinv}{\operatorname{dinv}}
\newcommand{\defc}{\operatorname{defc}}
\newcommand{\East}{\operatorname{East}}
\newcommand{\West}{\operatorname{West}}
\newcommand{\rev}{\operatorname{rev}}
\newcommand{\bk}{\operatorname{bk}}
\newcommand{\kap}{\mathcal K}

\title{West5 Nonfailure for Step-$\tau$ Dyck Words}
\author{}
\date{}

\begin{document}
\maketitle

\begin{abstract}
Fix an integer step $\tau\ge2$.  We prove that, for every length
$s\ge6$, the level-$5$ West move cannot fail when it is reached by the
successive leftmost-extraction procedure from a weak-lower-half
step-$\tau$ Dyck word of defect at most $\tau(s-2)-1$.  The proof uses
record payment, extraction transport, the far-pair inequalities from the
companion East5 argument, and a final split according to whether the word
before the third extraction begins with two zeros or with a positive
second entry.
\end{abstract}

\section{Definitions and imported facts}

Throughout, $\tau\ge2$ is fixed.  For words $U,V$, the notation $U:V$
means concatenation.

\begin{definition}[Normalized words and statistics]
A word $W=(w_0,\ldots,w_{n-1})$ is a \emph{normalized step-$\tau$ Dyck
word} if
\[
 w_0=0,\qquad w_i\ge0,\qquad w_{i+1}\le w_i+\tau
 \quad(0\le i<n-1).
\]
For nonnegative integers $u,v$, put
\[
 d_\tau(u,v)=
 \begin{cases}
  (u+\tau-v)_+,&u\le v,\\[1mm]
  (v+\tau+1-u)_+,&u>v,
 \end{cases}
\]
and
\[
 K_\tau(u,v)=\tau-d_\tau(u,v)
 =
 \begin{cases}
  \min(\tau,v-u),&u\le v,\\[1mm]
  \min(\tau,u-v-1),&u>v.
 \end{cases}
\]
For a word $W$ of length $n$, define
\[
 \area(W)=\sum_iw_i,
 \qquad
 \dinv_\tau(W)=\sum_{i<j}d_\tau(w_i,w_j),
\]
\[
 M_n=\tau\binom n2,
 \qquad
 \defc_\tau(W)=M_n-\area(W)-\dinv_\tau(W).
\]
Finally, set
\[
 \lambda_j(W)=\sum_{i<j}K_\tau(w_i,w_j)-w_j.
\]
Then every word satisfies
\begin{equation}
 \defc_\tau(W)=\sum_j\lambda_j(W).
 \label{eq:lambda-identity}
\end{equation}
\end{definition}

\begin{definition}[Extraction]
Let $W=(w_0,\ldots,w_{n-1})$ be normalized.  An occurrence
$w_j=e>0$ is \emph{extractable} if exactly one earlier occurrence lies
in
\[
 [\max(0,e-\tau),e),
\]
and deleting $w_j$ leaves a normalized word.  Equivalently, when
$0<j<n-1$, the deletion must satisfy the splice inequality
\[
 w_{j+1}\le w_{j-1}+\tau.
\]
At every stage we extract the leftmost extractable occurrence.
\end{definition}

\begin{definition}[The local tests]
The test $\East_3$ succeeds on $(u,v,w)$ exactly when
\[
 v\le w+\tau,
\]
and then leaves the triple unchanged.

Suppose $\East_3$ has failed on the inner triple of a five-window
$(a,b,c,p,q)$, so that $c>p+\tau$.  Put
\[
 \bk_\tau(b,c)=
 \begin{cases}
  (c,b),&b>c+\tau,\\
  (b,c),&b\le c+\tau.
 \end{cases}
\]
Then $\East_5$ is defined as follows.
\begin{enumerate}[label=\textup{(\alph*)}]
 \item If $b\le p+\tau$, it succeeds exactly when $b\le q+\tau$,
 and then sends
 \[
  (a,b,c,p,q)\longmapsto(a,p,c,b,q).
 \]
 \item If $b>p+\tau$, write $(b',c')=\bk_\tau(b,c)$.  It succeeds
 exactly when $c'\le q+\tau$, and then sends
 \[
  (a,b,c,p,q)\longmapsto(a,p,b',c',q).
 \]
\end{enumerate}
The West tests are obtained by reversal:
\[
 \West_j=\rev\circ\East_j\circ\rev
 \qquad(j=3,5).
\]
In particular, $\West_3$ succeeds on $(u,v,w)$ exactly when
$v\le u+\tau$.
\end{definition}

We use the following facts proved in the common preliminary work and in
the companion East5 proof.

\begin{lemma}[Imported record and extraction facts]
\label{lem:imported}
The following statements hold.
\begin{enumerate}[label=\textup{(\roman*)}]
 \item Let
 \[
 0=r_0<r_1<\cdots<r_m=M
 \]
 be the record values preceding a target $u$ in a normalized word.
 If $u\le M+\tau$, then
 \begin{equation}
  \sum_iK_\tau(r_i,u)\ge u+(M-1-u)_+.
  \label{eq:strong-record}
 \end{equation}
 In particular, every lambda increment of a normalized word is
 nonnegative.

 \item If $H\ge L+\tau+1$, then for every $u\ge0$,
 \begin{align}
  K_\tau(u,H)+K_\tau(u,L)
   &\ge \tau+\min\bigl(\tau,(L-u)_+\bigr),
   \label{eq:forward-far}\\
  K_\tau(H,u)+K_\tau(L,u)&\ge\tau.
   \label{eq:reverse-far}
 \end{align}

 \item If an extractable value $e$ is deleted and $e-1$ is appended,
 then
 \begin{equation}
  (\area,\dinv_\tau,\defc_\tau)
  \longmapsto
  (\area-1,\dinv_\tau+1,\defc_\tau).
  \label{eq:transport}
 \end{equation}
 These identities remain valid for the appended stage word even when it
 is not normalized.

 \item If $e_i,e_{i+1}$ are successive leftmost extracted values, then
 \begin{equation}
  e_i\le e_{i+1}+\tau.
  \label{eq:successive}
 \end{equation}

 \item A leftmost extractable occurrence is a record.  If a normalized
 word begins $(0,u,\ldots)$ with $u>0$, then all entries before its first
 extractable occurrence form a strictly increasing record chain and
 satisfy
 \begin{equation}
  w_{i+2}>w_i+\tau.
  \label{eq:rapid-propagation}
 \end{equation}

 \item If a normalized word of length $n$ begins with $(0,0)$ and its
 only extractable occurrence is the final one, then
 \begin{equation}
  \defc_\tau(W)\ge\tau(n-2).
  \label{eq:double-zero-boundary}
 \end{equation}
\end{enumerate}
\end{lemma}

\section{The common West5 estimates}

Suppose three successive leftmost extractions have values
\[
 x+1,\qquad y+1,\qquad z+1.
\]
After the third deletion, write the normalized remainder as
\[
 C=R:a:b,
 \qquad |R|=r.
\]
The corresponding level-$5$ stage is
\begin{equation}
 T=R:a:b:x:y:z.
 \label{eq:T-stage}
\end{equation}
In the application below, $r=s-5\ge1$.

\begin{lemma}[Consequences of West5 failure]
\label{lem:west-consequences}
Assume that $\West_3$ has failed before the third extraction and that
$\West_5$ fails on $(a,b,x,y,z)$.  Then
\begin{equation}
 x>b+\tau,
 \qquad x\le y+\tau,
 \qquad y\le z+\tau.
 \label{eq:history-ineqs}
\end{equation}
Moreover, exactly one of the following modes occurs:
\begin{align*}
 \mathrm{W1}:&\quad y\le b+\tau,\\[-1mm]
 &\quad y>a+\tau;\\[1mm]
 \mathrm{W2}:&\quad y>b+\tau,
 \quad y\le x+\tau,
 \quad x>a+\tau;\\[1mm]
 \mathrm{W3}:&\quad y>b+\tau,
 \quad y>x+\tau,
 \quad y>a+\tau.
\end{align*}
In modes W1 and W3, the far pairs are $(x,b)$ and $(y,a)$; in mode W2,
they are $(x,a)$ and $(y,b)$.  Thus each pair $(H,L)$ satisfies
$H\ge L+\tau+1$.

If
\[
\begin{aligned}
 U_W={}&K_\tau(a,b)+K_\tau(a,x)+K_\tau(a,y)\\
 &+K_\tau(b,x)+K_\tau(b,y)+K_\tau(x,y),
\end{aligned}
\]
then
\begin{equation}
 U_W\ge4\tau+K_\tau(a,b).
 \label{eq:UW}
\end{equation}
\end{lemma}

\begin{proof}
West3 failure gives $x>b+\tau$.  Equation
\eqref{eq:successive}, applied twice, gives the other two inequalities in
\eqref{eq:history-ineqs}.

West5 on $(a,b,x,y,z)$ is East5 on $(z,y,x,b,a)$.  Substitution into the
two East5 cases gives precisely W1--W3.

For \eqref{eq:UW}, consider the three modes.  In W1,
\[
 K_\tau(a,x)=K_\tau(a,y)=K_\tau(b,x)=\tau,
\]
and the reverse far-pair bound for $(x,b)$ at target $y$ gives
\[
 K_\tau(x,y)+K_\tau(b,y)\ge\tau.
\]
In W2,
\[
 K_\tau(a,x)=K_\tau(b,x)=K_\tau(b,y)=\tau,
\]
and the reverse far-pair bound for $(x,a)$ at target $y$ gives
\[
 K_\tau(x,y)+K_\tau(a,y)\ge\tau.
\]
In W3,
\[
 K_\tau(a,y)=K_\tau(b,x)=K_\tau(b,y)=K_\tau(x,y)=\tau.
\]
Adding the unused nonnegative term $K_\tau(a,b)$ proves
\eqref{eq:UW} in all cases.
\end{proof}

Define
\[
 E_a=\sum_{u\text{ in }R}\min\bigl(\tau,(a-u)_+\bigr),
 \qquad
 E_b=\sum_{u\text{ in }R}\min\bigl(\tau,(b-u)_+\bigr),
\]
\begin{equation}
 E=E_a+E_b,
 \qquad
 H=E+K_\tau(a,b)-a-b.
 \label{eq:E-H}
\end{equation}
Since $R:a:b$ is normalized, record payment gives
\begin{equation}
 E_a\ge a,
 \qquad
 E_b+K_\tau(a,b)\ge b,
 \qquad
 H\ge0.
 \label{eq:H-nonnegative}
\end{equation}
Also put
\begin{equation}
 \Delta=\defc_\tau(R),
 \qquad
 A_z=\sum_{u\text{ in }C}K_\tau(u,z)-z.
 \label{eq:Delta-Az}
\end{equation}
The unique predecessor of the extracted value $z+1$ survives in $C$.
Hence $z$ lies at most $\tau-1$ above some entry of $C$, and record
payment gives
\begin{equation}
 A_z\ge0.
 \label{eq:Az-nonnegative}
\end{equation}
The last lambda increment of $T$ is
\begin{equation}
 \lambda_z(T)=A_z+K_\tau(x,z)+K_\tau(y,z).
 \label{eq:lambda-z}
\end{equation}

\begin{lemma}[West suffix ledger]
\label{lem:west-ledger}
Let $\Lambda$ be the sum of the final five lambda increments of $T$.
If West5 fails, then
\begin{equation}
 \Lambda\ge
 2\tau r+4\tau+E+K_\tau(a,b)+\lambda_z(T)
 -(a+b+x+y).
 \label{eq:west-ledger}
\end{equation}
If, in addition,
\begin{equation}
 d:=\defc_\tau(T)\le\tau(r+3)-1,
 \label{eq:defect-cutoff-r}
\end{equation}
then
\begin{equation}
 x+y\ge
 \Delta+\tau(r+1)+1+\lambda_z(T)+H.
 \label{eq:xy-lower}
\end{equation}
\end{lemma}

\begin{proof}
For each occurrence $u$ in $R$, apply the forward far-pair inequality to
the two pairs supplied by Lemma~\ref{lem:west-consequences}.  This gives
\[
\begin{aligned}
 &K_\tau(u,a)+K_\tau(u,b)+K_\tau(u,x)+K_\tau(u,y)\\
 &\qquad\ge2\tau+
 \min\bigl(\tau,(a-u)_+\bigr)+
 \min\bigl(\tau,(b-u)_+\bigr).
\end{aligned}
\]
Summing over $R$, adding \eqref{eq:UW}, retaining the complete last
increment $\lambda_z(T)$, and subtracting $a+b+x+y$ proves
\eqref{eq:west-ledger}.

By \eqref{eq:lambda-identity}, $d=\Delta+\Lambda$.  Combine
\eqref{eq:west-ledger} with \eqref{eq:defect-cutoff-r}, then use the
definition of $H$ in \eqref{eq:E-H}, to obtain
\eqref{eq:xy-lower}.
\end{proof}

\begin{lemma}[Tail inequality]
\label{lem:tail}
If $x\le y+\tau$ and $y\le z+\tau$, then
\begin{equation}
 x+y\le
 2z+K_\tau(x,z)+K_\tau(y,z)+\tau+1.
 \label{eq:tail}
\end{equation}
Consequently, under the hypotheses of Lemma~\ref{lem:west-ledger},
\begin{equation}
 2z\ge\Delta+\tau r+A_z+H.
 \label{eq:z-master}
\end{equation}
In particular,
\begin{equation}
 z\ge\frac{\tau r}{2}.
 \label{eq:z-basic}
\end{equation}
\end{lemma}

\begin{proof}
Put $\phi_z(u)=u-K_\tau(u,z)$.  Directly from the definition of
$K_\tau$,
\[
 \phi_z(u)\le z\quad(u\le z),
 \qquad
 \phi_z(u)=z+1\quad(z<u\le z+\tau+1),
\]
and
\[
 \phi_z(u)=u-\tau\quad(u>z+\tau+1).
\]
Since $y\le z+\tau$, one has $\phi_z(y)\le z+1$.  If
$x\le z+\tau+1$, then
\[
 \phi_z(x)+\phi_z(y)\le2z+2\le2z+\tau+1.
\]
If $x>z+\tau+1$, then
\[
 \phi_z(x)=x-\tau\le y,
\]
so
\[
 \phi_z(x)+\phi_z(y)\le y+z+1\le2z+\tau+1.
\]
This is \eqref{eq:tail}.

Substitute \eqref{eq:lambda-z} into \eqref{eq:xy-lower} and compare with
\eqref{eq:tail}.  The terms $K_\tau(x,z)+K_\tau(y,z)$ cancel, giving
\eqref{eq:z-master}.  Equation \eqref{eq:z-basic} follows from
$\Delta,A_z,H\ge0$.
\end{proof}

For a word $W$, define its total chronological $K$-energy by
\begin{equation}
 \kap(W)=\sum_{i<j}K_\tau(w_i,w_j)
 =\area(W)+\defc_\tau(W).
 \label{eq:K-energy}
\end{equation}
For the stage \eqref{eq:T-stage}, put
\begin{equation}
 F=\kap(R)+E+K_\tau(a,b)+\lambda_z(T)+z.
 \label{eq:F}
\end{equation}

\begin{lemma}[Midpoint-energy reduction]
\label{lem:energy-reduction}
Assume \eqref{eq:defect-cutoff-r}.  Then
\begin{equation}
 J(T):=2\area(T)+\defc_\tau(T)
 \ge2F+3\tau r+5\tau+1.
 \label{eq:J-lower}
\end{equation}
If
\begin{equation}
 \Theta_r=\frac{\tau(r^2+3r+10)}4-\frac12,
 \label{eq:Theta}
\end{equation}
and $F>\Theta_r$, then
\begin{equation}
 J(T)>M_{r+5}.
 \label{eq:J-above}
\end{equation}
\end{lemma}

\begin{proof}
Rearrange \eqref{eq:west-ledger} and use
$\kap(R)=\area(R)+\Delta$ to obtain
\[
 \area(T)\ge F+2\tau r+4\tau-d,
 \qquad d=\defc_\tau(T).
\]
Thus
\[
 J(T)\ge2F+4\tau r+8\tau-d.
\]
Using $d\le\tau(r+3)-1$ gives \eqref{eq:J-lower}.
Finally,
\[
 2\Theta_r+3\tau r+5\tau+1
 =\tau\frac{(r+5)(r+4)}2=M_{r+5},
\]
which proves the last assertion.
\end{proof}

\section{The double-zero branch}

Let $D_2$ be the normalized word after the first two extractions and
before the third.  Write
\begin{equation}
 D_2=A:(z+1):B,
 \qquad |A|=h,
 \qquad |B|=k\ge1.
 \label{eq:D2-decomposition}
\end{equation}
The third extraction is nonfinal, so $B$ is nonempty.  Since
$D_2$ has length $s-2$,
\begin{equation}
 h+k=s-3.
 \label{eq:h-k}
\end{equation}

\begin{proposition}[Double-zero branch]
\label{prop:double-zero}
Suppose $D_2[1]=0$.  If West5 fails, then
\begin{equation}
 \defc_\tau(T)\ge\tau(s-2).
 \label{eq:double-zero-contradiction}
\end{equation}
\end{proposition}

\begin{proof}
The word $A:(z+1)$ begins with $(0,0)$ and has only its final occurrence
extractable.  Indeed, $z+1$ is the leftmost extractable occurrence of
$D_2$, and removing the suffix $B$ does not change the predecessor or
splice conditions of an earlier occurrence.  By
\eqref{eq:double-zero-boundary},
\[
 \defc_\tau(A:(z+1))\ge\tau(h-1).
\]
Transporting the final extraction gives
\begin{equation}
 \defc_\tau(A:z)\ge\tau(h-1).
 \label{eq:Az-base-defect}
\end{equation}
We prove
\begin{equation}
 \defc_\tau(T)-\defc_\tau(A:z)\ge\tau(k+2).
 \label{eq:double-zero-increment}
\end{equation}
Together with \eqref{eq:h-k} and \eqref{eq:Az-base-defect}, this yields
\eqref{eq:double-zero-contradiction}.

\medskip
\noindent\emph{Case 1: $k\ge2$.}
Write
\[
 B=B':a:b,
 \qquad
 P=A:B'.
\]
Then $|P|=r=h+k-2$.  Put
\[
 \mu=\defc_\tau(P)-\defc_\tau(A)\ge0;
\]
the inequality follows from lambda nonnegativity in the normalized word
$P$.

Apply the West suffix ledger with $P$ as prefix, and let $H_P\ge0$ be
the corresponding surplus from \eqref{eq:E-H}.  Subtracting
$\defc_\tau(A:z)$ gives
\[
\begin{aligned}
 \defc_\tau(T)-\defc_\tau(A:z)
 \ge{}&2\tau r+4\tau+\mu+H_P\\
 &+\sum_{v\text{ in }B}K_\tau(v,z)
 +K_\tau(x,z)+K_\tau(y,z)-x-y.
\end{aligned}
\]
By the tail inequality,
\begin{equation}
 \defc_\tau(T)-\defc_\tau(A:z)
 \ge2\tau r+3\tau-1+\mu+H_P
 +\sum_{v\text{ in }B}K_\tau(v,z)-2z.
 \label{eq:kge2-difference}
\end{equation}
Because $D_2$ begins with two zeros and $z+1$ occurs at index $h$,
\[
 z\le(h-1)\tau-1.
\]
Using $r=h+k-2$ in \eqref{eq:kge2-difference} and dropping nonnegative
terms gives
\[
\begin{aligned}
 \defc_\tau(T)-\defc_\tau(A:z)
 &\ge2\tau(h+k-2)+3\tau-1
   -2\bigl((h-1)\tau-1\bigr)\\
 &=2\tau k+\tau+1\\
 &\ge\tau(k+2).
\end{aligned}
\]
This proves \eqref{eq:double-zero-increment} when $k\ge2$.

\medskip
\noindent\emph{Case 2: $k=1$.}
Now $B=(b)$ and $A=R:a$.  Since $z+1$ is a record and its immediate
left neighbor $a$ is its unique predecessor, no earlier entry of $A$
can exceed $a$; hence $a$ is the record peak of $A$.  Moreover,
\begin{equation}
 a\le z\le a+\tau-1.
 \label{eq:a-z}
\end{equation}
The second initial zero gives
\begin{equation}
 a\le(h-2)\tau,
 \label{eq:a-height}
\end{equation}
and the splice condition for deleting $z+1$ gives
\begin{equation}
 b\le a+\tau.
 \label{eq:b-splice}
\end{equation}

For $t\ge0$, put
\[
 G_A(t)=\sum_{u\text{ in }A}K_\tau(u,t),
\]
\begin{equation}
 m_b=\min(\tau,b)+(a-1-b)_+,
 \label{eq:mb}
\end{equation}
and
\begin{equation}
 \gamma(t)=
 \begin{cases}
  t+\tau,&t\le a+\tau,\\
  h\tau,&t>a+\tau.
 \end{cases}
 \label{eq:gamma}
\end{equation}
Use the record chain of $A$ with the second initial zero omitted, then
add the contribution of that zero.  The strong record estimate gives
\begin{equation}
 G_A(b)-b\ge m_b.
 \label{eq:Gb}
\end{equation}
In every West5 failure mode, $x,y>\tau$.  Therefore the second zero
contributes a full $\tau$-unit to both targets, and
\begin{equation}
 G_A(t)\ge\gamma(t)
 \qquad(t=x,y).
 \label{eq:Ggamma}
\end{equation}

Set
\[
 I=\defc_\tau(T)-\defc_\tau(A:z).
\]
An exact lambda calculation gives
\[
\begin{aligned}
 I={}&[G_A(b)-b]+K_\tau(b,z)\\
 &+[G_A(x)+K_\tau(b,x)-x]\\
 &+[G_A(y)+K_\tau(b,y)+K_\tau(x,y)-y]\\
 &+K_\tau(x,z)+K_\tau(y,z).
\end{aligned}
\]
Hence, by \eqref{eq:Gb}--\eqref{eq:Ggamma},
\begin{equation}
\begin{aligned}
 I\ge\Psi:={}&m_b+K_\tau(b,z)+\gamma(x)-x+\gamma(y)-y\\
 &+K_\tau(b,x)+K_\tau(b,y)+K_\tau(x,y)\\
 &+K_\tau(x,z)+K_\tau(y,z).
\end{aligned}
\label{eq:Psi}
\end{equation}
We prove $\Psi\ge3\tau$ in the three failure modes.

\smallskip
\noindent\emph{Mode W1.}
Here $b>a$.  Put
\[
 \alpha=y-b,
 \qquad
 \beta=x-y,
 \qquad
 t=z-b.
\]
Then
\[
 1\le\alpha\le\tau,
 \qquad
 \tau-\alpha+1\le\beta\le\tau.
\]
Substitution in \eqref{eq:Psi} gives
\begin{equation}
 \Psi=m_b+K_\tau(b,z)+2(h\tau-b)+\tau-2
      +K_\tau(x,z)-t.
 \label{eq:Psi-W1}
\end{equation}
If $t\ge0$, then $K_\tau(b,z)=t$,
$K_\tau(x,z)\ge\tau-t$ by the reverse far-pair bound for $(x,b)$,
and
\[
 h\tau-b\ge\tau+1+t.
\]
Thus $\Psi\ge m_b+4\tau+t\ge3\tau$.  If $t<0$, then
\[
 K_\tau(b,z)=-t-1,
 \qquad
 K_\tau(x,z)=\tau,
 \qquad
 h\tau-b\ge\tau,
\]
so
\[
 \Psi\ge m_b+4\tau-3-2t\ge4\tau-1\ge3\tau.
\]

\smallskip
\noindent\emph{Mode W2 with $y\le a+\tau$.}
Put
\[
 u=x-a-\tau\ge1,
 \qquad
 v=a+\tau-y\ge0.
\]
Then $u+v\le\tau$ and
$K_\tau(x,y)=u+v-1$.  The pair $(y,b)$ is far, so
\[
 K_\tau(b,z)+K_\tau(y,z)\ge\tau.
\]
Moreover,
\[
 K_\tau(b,x)=K_\tau(b,y)=\tau.
\]
Writing $g=(h-2)\tau-a\ge0$, we have
\[
 \gamma(x)-x=\tau+g-u,
 \qquad
 \gamma(y)-y=\tau.
\]
Therefore
\[
 \Psi\ge m_b+5\tau+g+v-1+K_\tau(x,z)
 \ge5\tau-1\ge3\tau.
\]

\smallskip
\noindent\emph{Mode W2 with $y>a+\tau$.}
Put
\[
 u=x-a-\tau,
 \qquad
 v=y-a-\tau,
 \qquad
 \delta=z-a.
\]
Then
\[
 u,v\ge1,
 \qquad
 v\le\delta\le\tau-1,
 \qquad
 |u-v|\le\tau.
\]
Again
\[
 K_\tau(b,z)+K_\tau(y,z)\ge\tau,
 \qquad
 K_\tau(b,x)=K_\tau(b,y)=\tau,
\]
and
\[
 \gamma(x)-x=\tau+g-u,
 \qquad
 \gamma(y)-y=\tau+g-v.
\]
We claim
\begin{equation}
 K_\tau(x,y)+K_\tau(x,z)\ge u+v-\tau+1.
 \label{eq:ux-bound}
\end{equation}
If $u\le v$, the left side is
\[
 v-u+\tau+u-\delta-1,
\]
and \eqref{eq:ux-bound} follows from $u+\delta\le2\tau-2$.
If $v<u\le\delta$, the difference between the two sides of
\eqref{eq:ux-bound} is
\[
 2\tau+u-2v-\delta-3\ge0.
\]
If $u>\delta$, then the difference is
\[
 2(\tau-v-1)\ge0.
\]
Thus
\[
 \Psi\ge m_b+5\tau+2g-u-v+K_\tau(x,y)+K_\tau(x,z)
 \ge m_b+4\tau+2g+1\ge3\tau.
\]

\smallskip
\noindent\emph{Mode W3.}
If $x>a+\tau$, then
\[
 y>x+\tau>a+2\tau,
\]
contrary to $y\le z+\tau\le a+2\tau-1$.  Hence $x\le a+\tau$ and
$\gamma(x)-x=\tau$.  Also
\[
 \gamma(y)-y=h\tau-y\ge g+1.
\]
Furthermore,
\[
 K_\tau(b,x)=K_\tau(b,y)=K_\tau(x,y)=\tau,
\]
and the reverse far-pair bound for $(y,x)$ at target $z$ gives
\[
 K_\tau(x,z)+K_\tau(y,z)\ge\tau.
\]
Consequently
\[
 \Psi\ge5\tau+g+1\ge3\tau.
\]

Thus $I\ge3\tau=\tau(k+2)$ when $k=1$.  This proves
\eqref{eq:double-zero-increment} and completes the proposition.
\end{proof}

\section{The positive-second-entry branch}

\begin{lemma}[Rapid-prefix identities]
\label{lem:rapid-identities}
Assume $D_2[1]>0$ and write
\[
 D_2=A:(z+1):B,
 \qquad
 A=(q_0,\ldots,q_{h-1}),
 \qquad q_0=0.
\]
Put $q=q_{h-1}$.  Then $A:(z+1)$ is a strictly increasing rapid record
chain, and
\begin{equation}
 z+1-\tau\le q\le z,
 \qquad
 q_i\le z-\tau\quad(i\le h-2).
 \label{eq:q-z}
\end{equation}
Moreover,
\begin{align}
 \kap(A)&=q+\tau\binom{h-1}{2},
 \label{eq:KA}\\
 \sum_{u\text{ in }A}K_\tau(u,z)
 &=\tau(h-1)+z-q,
 \label{eq:A-to-z}\\
 \kap(A)+\sum_{u\text{ in }A}K_\tau(u,z)
 &=z+\tau\binom h2.
 \label{eq:rapid-energy}
\end{align}
For every $w\ge0$,
\begin{equation}
 \sum_{u\text{ in }A}K_\tau(u,w)+K_\tau(w,z)\ge z.
 \label{eq:rapid-target}
\end{equation}
\end{lemma}

\begin{proof}
Because $z+1$ is the first extractable occurrence and $D_2[1]>0$,
failed-splice propagation gives a strictly increasing record chain and
\eqref{eq:rapid-propagation}.  The immediate left neighbor $q$ lies in
the predecessor interval of $z+1$ and is its unique predecessor.  This
proves \eqref{eq:q-z}.

Only adjacent pairs of $A$ contribute less than a full $\tau$-unit to
$K$-energy.  The adjacent gaps telescope to $q$, while the nonadjacent
pairs contribute $\tau$.  This proves \eqref{eq:KA}.  By
\eqref{eq:q-z}, every occurrence before $q$ contributes $\tau$ to target
$z$, and $q$ contributes $z-q$, proving \eqref{eq:A-to-z}.  Their sum is
\eqref{eq:rapid-energy}.

For \eqref{eq:rapid-target}, first suppose $w\le q+\tau$.  Strong record
payment gives
\[
 \sum_{u\text{ in }A}K_\tau(u,w)
 \ge w+(q-1-w)_+.
\]
If $w\le q-1$, then
\[
 K_\tau(w,z)\ge z-q+1.
\]
If $q\le w\le z$, then $K_\tau(w,z)=z-w$.  If $w>z$, the prefix
contribution alone exceeds $z$.  These cases prove
\eqref{eq:rapid-target} when $w\le q+\tau$.

If $w>q+\tau$, every entry of $A$ contributes $\tau$, so the left side
is at least $h\tau$.  Normalization of $A:(z+1)$ gives
$z+1\le h\tau$, and \eqref{eq:rapid-target} follows.
\end{proof}

\begin{proposition}[Positive-second-entry branch]
\label{prop:positive}
Suppose $D_2[1]>0$, $r=s-5\ge1$, and West5 fails under the defect cutoff
\eqref{eq:defect-cutoff-r}.  Then
\begin{equation}
 F>\Theta_r.
 \label{eq:F-positive}
\end{equation}
\end{proposition}

\begin{proof}
Use the decomposition \eqref{eq:D2-decomposition}.

\medskip
\noindent\emph{Case 1: $k=|B|\ge2$ and $r\ge2$.}
Write
\[
 B=B':a:b,
 \qquad
 R=A:B'.
\]
Then $|B'|=r-h$.  Combine the exact identity
\eqref{eq:rapid-energy} with the target estimate
\eqref{eq:rapid-target}, once for each occurrence of $B'$.  Discard the
nonnegative internal $K$-energy of $B'$.  Finally, apply the two reverse
far-pair bounds at target $z$ and use
$E+K_\tau(a,b)\ge a+b$.  This gives
\begin{equation}
 F\ge(r-h+1)z+\tau\binom h2+2\tau+a+b.
 \label{eq:F-kge2}
\end{equation}

If $r\ge5$, substitute \eqref{eq:z-basic} into
\eqref{eq:F-kge2} and drop $a+b$.  Completing the square gives
\begin{equation}
\begin{aligned}
 F-\Theta_r
 \ge{}&\frac{\tau}{2}
 \left(h-\frac{r+1}{2}\right)^2\\
 &+\frac{\tau(r^2-4r-5)}8+\frac12>0.
\end{aligned}
\label{eq:large-r-square}
\end{equation}

It remains in this case to treat $2\le r\le4$.  Normalization gives
$z\le h\tau-1$, while \eqref{eq:z-basic} gives
$z\ge\tau r/2$.  Also $k\ge2$ implies $h\le r$.  Hence
\[
 (r,h)\in
 \{(2,2),(3,2),(3,3),(4,3),(4,4)\}.
\]
Using only $z\ge\tau r/2$ in \eqref{eq:F-kge2}, one obtains
\begin{equation}
 F-\Theta_r\ge a+b+
 \begin{cases}
  \frac12-\tau,&(r,h)=(2,2),(3,2),\\[1mm]
  \frac12-\frac\tau2,&(r,h)=(3,3),(4,3),\\[1mm]
  \frac12+\frac\tau2,&(r,h)=(4,4).
 \end{cases}
 \label{eq:small-r-table}
\end{equation}
Thus all these cases are finished if $a+b\ge\tau$.

Assume $a+b\le\tau-1$.  Since $z\ge\tau$, an elementary cap
calculation gives
\begin{equation}
 K_\tau(a,z)+K_\tau(b,z)\ge\min(z+1,2\tau).
 \label{eq:cap-ab}
\end{equation}
Indeed, if $z\ge2\tau$, both terms are $\tau$.  If
$\tau\le z\le2\tau-1$, either both terms are uncapped, in which case
their sum is at least $2z-(\tau-1)\ge z+1$, or a capped term only
increases the same lower bound.

By \eqref{eq:A-to-z},
\[
 A_z\ge\tau(h-1)-q+K_\tau(a,z)+K_\tau(b,z)
 \ge K_\tau(a,z)+K_\tau(b,z).
\]
Therefore \eqref{eq:z-master} and \eqref{eq:cap-ab} give
\[
 2z\ge\tau r+\min(z+1,2\tau).
\]
The alternative $z\le2\tau-1$ would imply
$z\ge\tau r+1\ge2\tau+1$, impossible.  Hence
\begin{equation}
 z\ge\frac{\tau(r+2)}2.
 \label{eq:z-improved}
\end{equation}
For $(r,h)=(2,2),(3,2),(4,3)$ this contradicts $z\le h\tau-1$.
For $(r,h)=(3,3)$, equation \eqref{eq:F-kge2} gives
\[
 F\ge z+5\tau\ge\frac{15\tau}{2}>7\tau-\frac12=\Theta_3.
\]
The pair $(4,4)$ was already positive in \eqref{eq:small-r-table}.
This completes Case~1.

\medskip
\noindent\emph{Case 2: $k=1$ and $r\ge2$.}
Here $A=R:a$.  Write
\[
 A=(q_0,\ldots,q_{r-1},a).
\]
The rapid-chain identities give
\begin{align}
 \kap(R)&=q_{r-1}+\tau\binom{r-1}{2},
 \label{eq:KR-k1}\\
 E_a&=\tau(r-1)+a-q_{r-1}.
 \label{eq:Ea-k1}
\end{align}
Every occurrence of $R$ lies outside the predecessor interval of
$z+1$, so
\begin{equation}
 \sum_{u\text{ in }R}K_\tau(u,z)=r\tau.
 \label{eq:R-to-z-k1}
\end{equation}
The two reverse far-pair bounds at target $z$ contribute another
$2\tau$, while $E_b+K_\tau(a,b)\ge b\ge0$.  Hence
\begin{equation}
 F\ge a+\tau\left(\binom{r-1}{2}+2r+1\right).
 \label{eq:F-k1}
\end{equation}
The unique-predecessor relation gives
\[
 a\ge z+1-\tau\ge\frac{\tau r}{2}+1-\tau.
\]
Substitution into \eqref{eq:F-k1} yields
\begin{equation}
 F-\Theta_r\ge
 \frac{\tau(r+3)(r-2)}4+\frac32>0.
 \label{eq:F-k1-positive}
\end{equation}

\medskip
\noindent\emph{Case 3: $r=1$.}
Now
\[
 R=(0),
 \qquad
 C=(0,a,b).
\]
Put
\begin{equation}
 H=\min(\tau,b)+K_\tau(a,b)-b\ge0.
 \label{eq:H-r1}
\end{equation}
Since $0\le a\le\tau$ and $b\le a+\tau$,
\begin{equation}
 a+b+H=
 \begin{cases}
  2a-1,&b<a,\\
  2b,&a\le b\le\tau,\\
  b+\tau,&b>\tau.
 \end{cases}
 \label{eq:abH}
\end{equation}
The two reverse far-pair bounds at target $z$ imply
\begin{equation}
 F\ge a+b+H+K_\tau(0,z)+2\tau.
 \label{eq:F-r1}
\end{equation}

There are two possible positions for the third extraction.

\smallskip
\noindent\emph{Subcase $h=1$.}
Then
\[
 D_2=(0,z+1,a,b),
 \qquad
 0\le z\le\tau-1.
\]
Here $K_\tau(0,z)=z$, and \eqref{eq:z-master} becomes
\begin{equation}
 2z\ge\tau+K_\tau(a,z)+K_\tau(b,z)+H.
 \label{eq:r1-h1-z}
\end{equation}
In particular, $z\ge\tau/2$.  We claim
\begin{equation}
 a+b+H+z\ge\frac{3\tau}{2}.
 \label{eq:r1-h1-claim}
\end{equation}
If $b<a$ and $a\le z$, equation \eqref{eq:r1-h1-z} reduces to
$2b+1\ge\tau$, hence $2a-1\ge\tau$.  If $b<a$ and $a>z$, then
$2a-1+z\ge3z+1$.  If $a\le b\le\tau$ and $b\le z$, equation
\eqref{eq:r1-h1-z} reduces to $2a\ge\tau$, hence $2b\ge\tau$.
If $a\le b\le\tau$ and $b>z$, then $2b+z\ge3z+2$.  Finally,
$b>\tau$ makes \eqref{eq:r1-h1-claim} immediate from
\eqref{eq:abH}.  Thus the claim holds.  Equations
\eqref{eq:F-r1} and \eqref{eq:r1-h1-claim} give
\[
 F\ge\frac{7\tau}{2}>\frac{7\tau}{2}-\frac12=\Theta_1.
\]

\smallskip
\noindent\emph{Subcase $h=2$.}
Then
\[
 D_2=(0,a,z+1,b),
 \qquad
 1\le a\le\tau,
 \qquad
 \tau\le z\le a+\tau-1.
\]
Now $K_\tau(0,z)=\tau$, and \eqref{eq:z-master} becomes
\begin{equation}
 2z\ge2\tau-a+K_\tau(b,z)+H.
 \label{eq:r1-h2-z}
\end{equation}
We claim
\begin{equation}
 a+b+H\ge\frac\tau2.
 \label{eq:r1-h2-claim}
\end{equation}
If $b<a$, then
$K_\tau(b,z)\ge z-a+1$.  Combining
\eqref{eq:r1-h2-z} with $z\le a+\tau-1$ gives
$2a+b\ge\tau+1$.  Since $b\le a-1$, this implies
$a\ge(\tau+2)/3$ and hence $2a-1\ge\tau/2$.  If
$a\le b\le\tau$, then $K_\tau(b,z)=z-b$, and
\eqref{eq:r1-h2-z} with $z\le a+\tau-1$ gives
$3a\ge\tau+1$; hence $2b\ge2a\ge\tau/2$.  If $b>\tau$,
\eqref{eq:r1-h2-claim} is immediate.  Therefore
\[
 F\ge3\tau+\frac\tau2
 =\frac{7\tau}{2}>\Theta_1.
\]

All cases give \eqref{eq:F-positive}.
\end{proof}

\section{West5 nonfailure}

\begin{theorem}[West5 nonfailure]
\label{thm:west5}
Let $\tau\ge2$ and $s\ge6$.  Let $D$ be a normalized step-$\tau$ Dyck
word satisfying
\begin{equation}
 d=\defc_\tau(D)\le\tau(s-2)-1,
 \qquad
 \area(D)\le\dinv_\tau(D).
 \label{eq:main-hypotheses}
\end{equation}
Suppose the DOWN procedure reaches its level-$5$ stage: the first two
leftmost extractions are made, West3 fails, and a third leftmost
extraction is made outside the final position.  Then the resulting
West5 test succeeds.
\end{theorem}

\begin{proof}
Suppose, for contradiction, that West5 fails.  Let the three extracted
values be $x+1,y+1,z+1$, and use the notation
\[
 T=R:a:b:x:y:z,
 \qquad
 r=s-5.
\]
By three applications of extraction transport,
\begin{equation}
 \area(T)=\area(D)-3,
 \qquad
 \dinv_\tau(T)=\dinv_\tau(D)+3,
 \qquad
 \defc_\tau(T)=d.
 \label{eq:three-transports}
\end{equation}
Since $\area(D)\le\dinv_\tau(D)$,
\[
 \dinv_\tau(T)-\area(T)
 =\dinv_\tau(D)-\area(D)+6\ge6.
\]
Equivalently,
\begin{equation}
 J(T)=2\area(T)+d\le M_s-6.
 \label{eq:J-below}
\end{equation}
Also $d\le\tau(r+3)-1$, so all the common estimates above apply.

Let $D_2$ be the word before the third extraction.  If $D_2[1]=0$,
Proposition~\ref{prop:double-zero} gives
\[
 d=\defc_\tau(T)\ge\tau(s-2),
\]
contrary to \eqref{eq:main-hypotheses}.

If $D_2[1]>0$, Proposition~\ref{prop:positive} gives
$F>\Theta_r$.  Lemma~\ref{lem:energy-reduction} then gives
\[
 J(T)>M_s,
\]
contrary to \eqref{eq:J-below}.

Both alternatives are impossible.  Therefore West5 succeeds.
\end{proof}

\end{document}
